% !PRINTED 669
\weberpage{669}
\begin{center}
§ 190. \gesperrt{Branch Points and Branch Numbers.}
\end{center}

According to § 184 there is only a finite number of points at which $z-c$ becomes zero in an order higher than the first. These are called the branch points of $\Omega$ in $z$.

We construct a polygon $\mathfrak{Z}_z$ which contains $(e-1)$ times every point $\mathfrak{P}$ at which $z-c$ or $1:z$ becomes zero in the $e$th order. This polygon is called the \gesperrt{branch polygon of $\Omega$ with respect to $z$}, and its order

\weq{2}{w_z = \Sigma(e-1)}

is called the \gesperrt{branch number with respect to $z$}\footnote{If one spreads the \gesperrt{Riemann} surface over the $z$-plane, then $n$ is the number of sheets lying one above another, and $w_z$ is the number of branch points of this surface.}.

In § 184, (20) we obtained the formula:

\weq{3}{N(t-\tau) \equiv (t-\tau_1)(t-\tau_2)\ldots(t-\tau_n) \ (\mathrm{mod}\ z-c),}

in which $\tau$ denotes the basis form of the field $\Omega$ with respect to $z$, and $\tau_1, \tau_2, \ldots, \tau_n$ the equal or different constant functionals to which $\tau$ becomes congruent with respect to the prime factors of $z-c$, that is to say, the values which $\tau$ takes at the points $\mathfrak{P}_1, \mathfrak{P}_2, \ldots, \mathfrak{P}_n$ conjugate with respect to $z$. At each of these points (3) must pass over into a correct equation, and from this there follow, for example:

\weq{4}{N(\tau) = \tau_1 \tau_2 \ldots \tau_n \,,}

\weq{5}{S(\tau) = \tau_1 + \tau_2 + \cdots + \tau_n}

as the values of the rational functionals $N(\tau)$ and $S(\tau)$ for $z=c$. If one substitutes rational functions of $z$ for the functional variables, then corresponding formulae hold for the functions in $\Omega$, provided that none of the exceptional cases $0\,.\,\infty$, $\infty+\infty$ occurs.

If, therefore, $\eta_1, \eta_2, \ldots, \eta_n$ is a basis of $\Omega$, and if $\eta_{i,1}, \eta_{i,2}, \ldots, \eta_{i,n}$ are values of $\eta_i$ conjugate with respect to $z$, then by (5) and by the definition of the discriminant in § 173

\weq{6}{\Delta(\eta_1, \eta_2, \ldots, \eta_n) = (\Sigma \pm \eta_{1,1}\eta_{2,2}\ldots\eta_{n,n})^2.}

If two functions $z$ and $z'$ in $\Omega$ stand in a linear relation:

\weq{7}{z' = \dfrac{az+b}{cz+d},}
%CONT>>
